(1) 对连续性求时间导数得 \(c_vT_{tt}+q_{xt}=0\)，并由Cattaneo式 \(q_t=(-q-\kappa T_x)/\tau\) 消去 \(q_{xt}\)。利用 \(q_x=-c_vT_t\)，得到
\[
\boxed{\tau T_{tt}+T_t=D T_{xx}},\quad D=\kappa/c_v.
\]
此为telegraph双曲方程，必须指定 \(T(x,0)\) 与 \(T_t(x,0)\)（等价于相容初始热流散度），而Fourier热扩散只需初温；热波速度有限依赖模型尺度，非任意频段的真实晶格声子速度。

(2) 代入 \(e^{i(kx-\omega t)}\) 得 \(-\tau\omega^2-i\omega+Dk^2=0\)，故
\[
\omega_\pm=\frac{-i\pm\sqrt{4\tau Dk^2-1}}{2\tau}.
\]
当 \(4\tau Dk^2>1\) 时两个模式传播的实部为 \(\pm\sqrt{4\tau Dk^2-1}/(2\tau)\)，高波数相速度渐近 \(\pm\sqrt{D/\tau}\)，衰减率 \(1/(2\tau)\)。长波展开则为扩散慢模 \(-iDk^2+O(k^4)\) 与迅速弛豫模 \(-i/\tau+iDk^2+O(k^4)\)，不能把两个分支都当作无损热声波。

(3) 将 \(c_vT_t+q_x=0\) 乘以 \(T\)，将 \(\tau q_t+q=-\kappa T_x\) 乘以 \(q/\kappa\)，相加并整理：
\[
\boxed{\partial_t\left(\frac{c_vT^2}{2}+\frac{\tau q^2}{2\kappa}\right)
+\partial_x(Tq)=-\frac{q^2}{\kappa}\le0}.
\]
周期边界或空间无穷远快速衰减时，总积分单调下降。该二次量是线性化温度扰动的Lyapunov泛函而非温度正负任意时均为物理内能 \(c_vT\)；将二者混同会导致维度和符号误判。

(4) 固定 \(D,k\) 令 \(\tau\downarrow0\)，慢根保持 \(\omega_{\rm slow}\to-iDk^2\)，快根虚部 \(-1/\tau\) 趋负无穷，热流迅速逼近 \(q=-\kappa T_x\)，恢复Fourier扩散。但若先取 \(k\to\infty\) 则获得具有 \(v=\sqrt{D/\tau}\) 的传播分支，接着令 \(\tau\to0\) 会使模型传播速度发散。故该极限非一致，不能在所有短波范围上把无穷速Fourier方程当成有限速热波的精确极限。
